\documentclass[a4paper,12pt]{article}
\usepackage[norwegian]{babel}
\usepackage{amsmath, amssymb}
\usepackage{geometry}
\usepackage{fancyhdr}
\usepackage{listings}
\usepackage{xcolor}
\usepackage{enumitem}
\newcommand{\D}{\mathrm{d}}
\usepackage[colorlinks=true, linkcolor=red!50!black, citecolor=red!50!black, urlcolor=blue]{hyperref}

\geometry{margin=2.5cm}
\pagestyle{fancy}
\fancyhf{}
\rhead{R1}
\lhead{Parameterframstilling}
\rfoot{\thepage}

\usepackage{setspace}
\onehalfspacing

\lstset{
    basicstyle=\ttfamily\small,
    backgroundcolor=\color{gray!10},
    frame=single,
    breaklines=true,
    columns=fullflexible,
    keepspaces=true,
    showstringspaces=false,
    upquote=true
}

\title{Introduksjon til parameterframstilling og vektorfunksjoner}
\author{}
\date{}

\begin{document}

\maketitle


\section*{Newtons første lov}

I én dimensjon er rettlinjet bevegelse med konstant fart beskrevet av en lineær funksjon som 
\[
x(t) = 2t-2,
\]
der $x$ er posisjon i meter og $t$ er tid i sekunder. Vi finner farta $v$ ved å se på stigningstallet eller ved å derivere. Uansett får vi at
\[
v(t) = \frac{dx}{dt} = 2,
\]
og farta er med andre ord konstant 2 m/s. Akselerasjonen $a$ finner vi ved å derivere farta, hvilket gir
\[
a(t) = \frac{dv}{dt} = 0.
\]
Akselerasjonen er altså konstant 0 m/s$^2$.

Om akselerasjonen er lik null, vil et objekt bevege seg i rett linje med konstant fart. Dette kalles ofte Newtons første lov.

\section*{To dimensjoner}

Du skal nå se på rettlinjet bevegelse med konstant fart i to dimensjoner.

\subsection*{Oppgave 1: Parameterframstilling}

Det er flere måter å beskrive bevegelse i to dimensjoner. Den første vi skal se på er \textit{parameterframstilling}. Denne minner også mest om eksemplet fra én dimensjon. La

\[
l:
\begin{cases}
x(t)=2t - 2 \\
y(t)=\phantom{2}t-2
\end{cases},\, 0 \leq t \leq 10.
\]
Vi sier at linja $l$ er gitt ved $x(t)$ og $y(t)$ for $t$ mellom 0 og 10.

\begin{enumerate}[label=\alph*)]
\item For å bruke parameterframstilling i GeoGebra må du bruke denne funksjonen:
\begin{lstlisting}
Kurve(<Uttrykk>,<Uttrykk>,<Parametervariabel>,<Startverdi>,<Sluttverdi>)
\end{lstlisting}
Skriv følgende i GeoGebra:
\begin{lstlisting}
l = Kurve(2t-2, t-2, t, 0, 10)
\end{lstlisting}

\item Hva representerer endepunktene på grafen til $l$?
\end{enumerate}


\subsection*{Oppgave 2: Punkt}

Vi kan også beskrive bevegelsen som et punkt $A$ med koordinater $(x(t),y(t))$. I dette tilfellet blir koordinatene $(2t-2,t-2)$.

\begin{enumerate}[label=\alph*)]
\item Legg til en \lstinline{Glider} for $t$. Skriv følgende i GeoGebra:
\begin{lstlisting}
t = 0
\end{lstlisting}
Høyereklikk på \lstinline{t} og velg \lstinline{Lag glider}.

Høyereklikk på \lstinline{t} igjen og velg \lstinline{Innstillinger}. Gå til \lstinline{Glider} og sett \lstinline{Min} til 0, \lstinline{Maks} til 10, og \lstinline{Animasjonstrinn} til 0.01.

\item Legg til punktet $A$, skriv:
\begin{lstlisting}
A = (l(t))
\end{lstlisting}
Dette har samme effekt som å skrive \lstinline{A = (2t-2,t-2)}.
\item Animer $t$ ved å trykke på startknappen. Hva slags bevegelse observerer du?
\item Hvor langt har punktet $A$ beveget seg fra $t=0$ til $t=10$?\footnote{$\Delta s = 10\sqrt{5}$ m $\approx 22.4$ m.}
\item Hva er hastigheten?\footnote{$v = \sqrt 5$ m/s $\approx 2.24$ m/s.}
\end{enumerate}

\subsection*{Oppgave 3: Vektorfunksjon}

Den siste\footnote{Men kanskje også beste måten.} måten å beskrive bevegelsen på er med en \textit{vektorfunksjon}. Denne får navnet posisjonsvektor $\vec r(t)$, og er gitt ved
\[\vec r(t) = \overrightarrow{OA}.\]
Med andre ord er $\vec r (t)$ en vektor som peker på punktet der det er akkurat nå. Vi bruker vanlig vektornotasjon:
\[
\vec r(t) = 
\begin{pmatrix}
x(t) \\
y(t)
\end{pmatrix} =
\begin{pmatrix}
2t-2 \\
t-2
\end{pmatrix}.
\]

\begin{enumerate}[label=\alph*)]
\item Legg til $\vec r (t)$ i GeoGebra. Skriv følgende:
\begin{lstlisting}
r(t) = Vektor(A)
\end{lstlisting}

\item Fartsvektoren $\vec v(t)$ er gitt som den deriverte av posisjonsvektoren $\vec r(t)$:
\[ \vec v(t) = \frac{d\vec r}{dt}. 
\]
Vis / argumenter for / velg å tro at \[ \vec v(t) = \begin{pmatrix}
2 \\
1
\end{pmatrix}.\]

\item Det går selvsagt også an å gjøre dette i GeoGebra. Skriv følgende:
\begin{lstlisting}
b = Derivert(l)
v = Vektor(A, A + b(t))
\end{lstlisting}
Det kan være en idé å endre farge på $\vec v$. Det er også lurt å skjule $b(t)$ siden denne kun er en hjelpefunksjon for å nå $\vec v$.\newline

Finn hastigheten $|\vec v(t)|$ og sammenlign med 2 e).

\item Bruk samme strategi for akselerasjonsvektoren $\vec a(t)$. Altså skriv:
\begin{lstlisting}
c = Derivert(b)
a = Vektor(A, A + c(t))
\end{lstlisting}

Hvorfor tror du $\vec a$ ikke synes?



\end{enumerate}


\section*{Bølgebevegelse}

Siden du har vært flink pike/dreng og gjort som fortalt, kan du nå enkelt studere andre bevegelser. La f.eks.

\[
l:
\begin{cases}
x(t)=t \\
y(t)=\sin(t)
\end{cases},\, 0\leq t \leq 10.
\]


\begin{enumerate}[label=\alph*)]
\item Alt du må gjøre er å endre på definisjonen av $l$ slik at det står:
\begin{lstlisting}
l = Kurve(t, sin(t), t, 0, 10)
\end{lstlisting}
i stedet for:
\begin{lstlisting}
l = Kurve(2t-2, t-2, t, 0, 10)
\end{lstlisting}

\item Hva blir utrykket for posisjonsvektoren $\vec r(t)$?
\item Hva blir utrykket for fartsvektoren $\vec v(t)$ og akselerasjonsvektoren $\vec a(t)$?

\item Når er akselerasjonen størst? Hvordan ser grafen ut da? Hva har dette med bilkjøring å gjøre?

\end{enumerate}

Det er nok ikke dumt å lagre denne GeoGebra-fila.

\section*{Båtkollisjon}
Du skal i det som følger finne en måte å avverge båtkollisjoner på. Du burde åpne et nytt GeoGebra-vidu.

\subsection*{1. Tilfelle}
To båter segler med konstant fart på havet. Båtene er gitt som $a$ og $b$ og de har parameterframstilling:

\[
a:
\begin{cases}
x(t)=2t - 2 \\
y(t)=\phantom{2}t-2
\end{cases},\, 0 \leq t \leq 10; \,\,\,\,\,\,\,\,\,\,\,\,\,\,\,\,\, b:
\begin{cases}
x(t)=-t + 10 \\
y(t)=\,\,3t-4
\end{cases},\, 0 \leq t \leq 10
\]

\begin{enumerate}[label=\alph*)]
\item Legg både $a$ og $b$ inn i GeoGebra. 
\begin{lstlisting}
a = Kurve(2t-2, t-2, t, 0, 10)
b = Kurve(-t+10, 3t-4, t, 0, 10)
\end{lstlisting}
Tror du båtene vil kollidere?
\item Legg til en glider for $t$. Legg deretter til punktene $A$ og $B$.
\begin{lstlisting}
A = (a(t))
B = (b(t))
\end{lstlisting}
Kolliderer båtene?
\item Finn $\overrightarrow{AB}$. I GeoGebra kan du skrive:
\begin{lstlisting}
AB = Vektor(A,B)
\end{lstlisting}
Hva representerer denne vektoren?

\item Vis ved regning at $\overrightarrow{AB}$ er gitt ved
\[
\overrightarrow{AB} =
\begin{pmatrix}
-3t + 12 \\
2t -2
\end{pmatrix}.
\]
Hvordan kan du bruke dette til å vise at båtene ikke kolliderer?
\end{enumerate}


\subsection*{2. Tilfelle}
To båter segler med konstant fart på havet. Båtenes bevegelser er gitt ved parameterframstilling:

\[
l:
\begin{cases}
x(t)=3t - 2 \\
y(t)=\phantom{3}t-6
\end{cases},\, 0 \leq t \leq 10; \,\,\,\,\,\,\,\,\,\,\,\,\,\,\,\,\, m:
\begin{cases}
x(t)=2t + 5 \\
y(t)=-8t+57
\end{cases},\, 0 \leq t \leq 10
\]

\begin{enumerate}[label=\alph*)]
\item Finn ut om båtene kolliderer.
\item Legg inn vektoren $\vec r(t)$ mellom båtene.
\item Vis at denne vektoren er gitt ved:
\[
\vec r(t) =
\begin{pmatrix}
-t+7 \\
-9t +63
\end{pmatrix} = (t-7)
\begin{pmatrix}
-1\\
-9
\end{pmatrix}.
\]
Hvordan kan du se at disse to båtene vil kollidere?

\item Bruk det du har lært til å foreslå hvordan båter kan finne ut om de er på kollisjonskurs, og også foreslå hvordan de kan unngå å kollidere når de har funnet ut at de er på en slik kollisjonskurs.

\end{enumerate}

\section*{Fra punkt til linje}

Ei vanlig oppgave er å måtte finne avstanden fra punkt til linje.

\subsection*{Oppgave 1:}

En båt $B$ segler med konstant fart på havet. Posisjonen er gitt ved parameterframstillinga
\[
b:
\begin{cases}
x(t)=3t \\
y(t)=-8t+5
\end{cases}.
\]

\begin{enumerate}[label=\alph*)]
\item Legg inn i GeoGebra parameterframstillinga $b$, punktet $B$, og posisjonsvektoren $\vec r(t) = \overrightarrow{OB}$. For punktet og vektoren må du huske å legge inn en glider for $t$.

\item Bruk posisjonsvektoren 
\[
\vec r(t)= 
\begin{pmatrix}
3t\\
-8t+5
\end{pmatrix}
\]

til å vise at avstanden fra origo $O$ til båten $B$ er gitt ved
\[
|\vec r(t)| = \sqrt{73t^2-80t +25}.
\]

\item Når er båten nærmest origo?\footnote{Hint: Her må du derivere for å finne ut når $|\vec r(t)|$ er minst. Avstanden er minst når også den kvadrerte avstanden $|\vec r(t)|^2$ er minst. Altså kan du derivere $|\vec r(t)|^2$ i stedet for $|\vec r(t)|$ for å løse for $t$.}\textsuperscript{,}\footnote{$t_{\text{min}}= \frac{10}{19} \approx 0.526$.} Hvor langt unna er den da?\footnote{$|\vec r(t_{\text{min}})| = \frac{15\sqrt 5}{19} \approx 1.77$.}

Gitt punkt $A(3,4)$.

\item Vis at vektoren fra punktet til båten er gitt ved \[
\overrightarrow{AB} = 
\begin{pmatrix}
3t-3\\
-8t+1
\end{pmatrix}
\]
\item Hvor nærme kommer båten $A$?\footnote{$|\overrightarrow{AB}_{\text{min}}| = \frac{21\sqrt{73}}{73} \approx 2.46.$}

\item Hvor nærme hverandre kommer de to båtene fra 1. tilfelle over?\footnote{$|\vec s_\text{min}| = \frac{18\sqrt{13}}{13} \approx 4.99.$} Husk at separasjonsvektoren mellom båtene $\vec s$ var gitt ved
\[
\vec s =
\begin{pmatrix}
-3t + 12 \\
2t -2
\end{pmatrix}.
\]

\end{enumerate}
\newpage
\subsection*{Oppgave 2:}

Gitt linja
\[
y = -2x+10.
\]

\begin{enumerate}[label=\alph*)]
\item Finn ei parameterframstilling for linja.\footnote{Du kan velge denne fritt. Et enkelt valg er $x(t) = t$, da blir $y(t)=-2t +10$.}
\item Finn et utrykk for posisjonsvektoren $\vec r(t)$ for denne parametriseringa.
\item Finn korteste avstand fra origo til linja.\footnote{$|\vec r_\text{min}| = 2\sqrt 5 \approx 4.47.$}

\item Finn korteste avstand fra punktet $(5,0)$ til linja.
\end{enumerate}
\end{document}